W ciągu ostatnich lat sztuczna inteligencja dynamicznie przkształciła wiele obszarów gospodarki i życia codziennego. Systemy oparte na sieciach neuronowych wspierają wyszukiwarki internetowe, umożliwiają autonomiczną jazdę samochodów i optymalizują procesy produkcyjne. Rosnąca zdolność tych modeli do analizowania ogromnych zbiorów danych oraz uczenia się z nich pozwala na podejmowanie skomplikowanych decyzji bez bezpośredniego nadzoru człowieka. Przykładowo, w sektorze finansowym sieci neuronowe wspomagają zarządzanie portfelem inwestycyjnym, prognozowanie ryzyka kredytowego czy dynamiczne kształtowanie cen. Wyzwanie polega jednak na tym, by w środowiskach o wysokim stopniu niepewności i dużej liczbie dostępnych opcji sieć potrafiła skutecznie eksplorować przestrzeń możliwych działań i wypracować stabilną, optymalną strategię działania.

Gry stanowią idealną platformę testową do weryfikacji i porównywania metod uczenia maszynowego. Pierwszym wyzwaniem, któremu podołała sztuczna inteligencja były szaczy. W latach roku 2015 AlphaGo, został pierwszym programem, który pokonał profesjonalnego gracza go bez żadnego wsparcia. 
